%% kapa: exam=1 round=33 year=2022 period=2 form=A subject=accounting subject_ko=회계학 seq=13 number=53
%% source: Q-Net 감정평가사 기출문제 (cst00309) artlSeq=5209997 file=raw/1008/5209997/2022년도 제33회 감정평가사 1차시험 2교시(감정평가 관계법규, 회계학).pdf
%% conversion: kordoc 4.24.51 -> tools/gen1.py (auto); review=proofread-sonnet; flags=none
\begin{question}{53}
(주)감평(리스이용자)은 20×1년 1월 1일에 (주)한국리스(리스제공자)와 다음과 같은 리스계약을 체결하였다.
\begin{framedtext}
○ 리스개시일: 20×1년 1월 1일 \par
○ 리스기간: 20×1년 1월 1일부터 20×3년 12월 31일까지 \par
○ 고정리스료: 매년 말 ￦1,000,000 후급 \par
○ (주)감평은 리스기간 종료일에 (주)한국리스에게 ￦300,000을 지급하고, 기초자산(리스자산)의 소유권을 이전 받기로 하였다. \par
○ (주)감평과 (주)한국리스는 리스개시일에 리스개설직접원가로 각각 ￦100,000과 ￦120,000을 지출하였다. \par
○ 리스개시일 현재 기초자산의 내용연수는 4년이고, 잔존가치는 ￦0이다. \par
\end{framedtext}
(주)감평은 사용권자산에 대해 원가모형을 적용하고 있으며 정액법으로 감가상각한다. 리스 관련 내재이자율은 알 수 없으나 (주)감평의 증분차입이자율이 연 10\%라고 할 때, 상기 리스거래와 관련하여 (주)감평이 20×1년도에 인식할 비용총액은? (단, 상기 리스계약은 소액 기초자산 리스에 해당하지 않으며, 감가상각비의 자본화는 고려하지 않는다. 또한, 단수차이로 인한 오차가 있다면 가장 근사치를 선택한다.)
\par{\small\centering\begin{tabular}{|c|c|c|}\hline
\multirow{2}{*}{기간} & 단일금액 ￦1의 현재가치 & 정상연금 ￦1의 현재가치 \\ \cline{2-3}
 & 10\% & 10\% \\ \hline
3 & 0.75131 & 2.48685 \\ \hline
\end{tabular}\par}
\begin{choices}
  \item ￦532,449
  \item ￦949,285
  \item ￦974,285
  \item ￦1,175,305
  \item ￦1,208,638
\end{choices}
\end{question}
\answer{3}
